% GENERATED FILE. Edit canonical lessons, then run npm run generate:books.
% canonical-source-hash: fnv1a64:72f412d5a4f418f5
% canonical-lessons: HI-C337-khazana, HI-C337-garna, HI-C337-nag, HI-C337-mani, HI-C337-dasna

\chapter{Buried Treasure}
\label{ch:hi-buried-treasure}
\hlchaptermodality{337}

\begin{chapteropening}
\textbf{By the end of this chapter:} \emph{I can tell a tale of hidden treasure.}

Complete the last lesson of chapter 337: I can tell a tale of hidden treasure.
\end{chapteropening}

\section[\texorpdfstring{\hi{ख़ज़ाना} (khazānā)}{ख़ज़ाना (khazānā)}]{\hi{ख़ज़ाना} (khazānā) — a treasure; a treasury}

\label{lesson:HI-C337-khazana}



\begin{quote}
Before the new one: say the Hindi for to regret, then the Hindi for to sulk.
\end{quote}

\subsection*{You'll want to know: \hi{ख़ज़ाना}}
\textbf{\hi{ख़ज़ाना}} — \emph{khazānā} — \textquotedblleft{}a treasure; a treasury\textquotedblright{}.

\begin{quote}
\hi{कहानी} \hi{में} \hi{राजा} \hi{का} \hi{ख़ज़ाना} \hi{एक} \hi{पुराने} \hi{क़िले} \hi{में} \hi{था।}
\end{quote}

\begin{quote}
\emph{kahānī meṁ rājā kā khazānā ek purāne qile meṁ thā.}
\end{quote}

\begin{quote}
In the story the king's treasure was in an old fort.
\end{quote}

Masculine.

\begin{cousinweb}
\textbf{\hi{ख़ज़ाना}} (\emph{khazānā}) is usually traced, through Persian, to Arabic \emph{khizāna}, \textquotedblleft{}a storehouse\textquotedblright{}. It is a king's treasury and, in stories, gold hidden underground.
\end{cousinweb}

\begin{grammarlens}[title={What you've built}]
One more word for telling what happened.
\end{grammarlens}

\subsection*{Your turn}
\begin{itemize}
\raggedright
  \item \emph{Say it:} \emph{khazānā}
  \item \emph{Say it:} \emph{khazānā}, once more
  \item \emph{Say it:} \emph{kahānī meṁ rājā kā khazānā ek purāne qile meṁ thā.} — In the story the king's treasure was in an old fort.
  \item \emph{Recall:} say the Hindi for to regret, then the Hindi for to sulk, then say \emph{khazānā} again
\end{itemize}

\begin{tcolorbox}[breakable,colback=teal!4,colframe=teal!35!black,title={Before you move on}]
What does \hi{ख़ज़ाना} mean? (\textquotedblleft{}A treasure; a treasury\textquotedblright{}.) Say it once more.
\end{tcolorbox}
\section[\texorpdfstring{\hi{गाड़ना} (gāṛnā)}{गाड़ना (gāṛnā)}]{\hi{गाड़ना} (gāṛnā) — to bury}

\label{lesson:HI-C337-garna}



\begin{quote}
Before the new one: say the Hindi for to sulk, then the Hindi for a treasure.
\end{quote}

\subsection*{You'll want to know: \hi{गाड़ना}}
\textbf{\hi{गाड़ना}} — \emph{gāṛnā} — \textquotedblleft{}to bury\textquotedblright{}.

\begin{quote}
\hi{राजा} \hi{ने} \hi{ख़ज़ाना} \hi{एक} \hi{पेड़} \hi{के} \hi{नीचे} \hi{गाड़ा।}
\end{quote}

\begin{quote}
\emph{rājā ne khazānā ek peṛ ke nīce gāṛā.}
\end{quote}

\begin{quote}
The king buried the treasure under a tree.
\end{quote}

Past, with \hi{ने} (he or she …): \hi{उसने} \hi{पैसा} \hi{गाड़ा}, \emph{usne paisā gāṛā}.

\begin{cousinweb}
\textbf{\hi{गाड़ना}} (\emph{gāṛnā}) is inherited from \emph{*gaḍḍ}, a form CDIAL reconstructs (CDIAL 3979): to drive something into the ground, a post or a treasure.
\end{cousinweb}

\begin{grammarlens}[title={What you've built}]
One more word for telling what happened.
\end{grammarlens}

\subsection*{Your turn}
\begin{itemize}
\raggedright
  \item \emph{Say it:} \emph{gāṛnā}
  \item \emph{Say it:} \emph{gāṛnā}, once more
  \item \emph{Say it:} \emph{rājā ne khazānā ek peṛ ke nīce gāṛā.} — The king buried the treasure under a tree.
  \item \emph{Recall:} say the Hindi for to sulk, then the Hindi for a treasure, then say \emph{gāṛnā} again
\end{itemize}

\begin{tcolorbox}[breakable,colback=teal!4,colframe=teal!35!black,title={Before you move on}]
What does \hi{गाड़ना} mean? (\textquotedblleft{}To bury\textquotedblright{}.) Say it once more.
\end{tcolorbox}
\section[\texorpdfstring{\hi{नाग} (nāg)}{नाग (nāg)}]{\hi{नाग} (nāg) — a cobra; a serpent}

\label{lesson:HI-C337-nag}



\begin{quote}
Before the new one: say the Hindi for a treasure, then the Hindi for to bury.
\end{quote}

\subsection*{You'll want to know: \hi{नाग}}
\textbf{\hi{नाग}} — \emph{nāg} — \textquotedblleft{}a cobra; a serpent\textquotedblright{}.

\begin{quote}
\hi{ख़ज़ाने} \hi{के} \hi{पास} \hi{एक} \hi{नाग} \hi{रहता} \hi{था।}
\end{quote}

\begin{quote}
\emph{khazāne ke pās ek nāg rahtā thā.}
\end{quote}

\begin{quote}
A cobra lived beside the treasure.
\end{quote}

Masculine.

\begin{cousinweb}
is Sanskrit \emph{nāga}, \textquotedblleft{}a serpent, a cobra\textquotedblright{} (Monier-Williams). In tales the \textbf{\hi{नाग}} guards hidden treasure.
\end{cousinweb}

\begin{grammarlens}[title={What you've built}]
One more word for telling what happened.
\end{grammarlens}

\subsection*{Your turn}
\begin{itemize}
\raggedright
  \item \emph{Say it:} \emph{nāg}
  \item \emph{Say it:} \emph{nāg}, once more
  \item \emph{Say it:} \emph{khazāne ke pās ek nāg rahtā thā.} — A cobra lived beside the treasure.
  \item \emph{Recall:} say the Hindi for a treasure, then the Hindi for to bury, then say \emph{nāg} again
\end{itemize}

\begin{tcolorbox}[breakable,colback=teal!4,colframe=teal!35!black,title={Before you move on}]
What does \hi{नाग} mean? (\textquotedblleft{}A cobra; a serpent\textquotedblright{}.) Say it once more.
\end{tcolorbox}
\section[\texorpdfstring{\hi{मणि} (maṇi)}{मणि (maṇi)}]{\hi{मणि} (maṇi) — a jewel; a gem}

\label{lesson:HI-C337-mani}



\begin{quote}
Before the new one: say the Hindi for to bury, then the Hindi for a cobra.
\end{quote}

\subsection*{You'll want to know: \hi{मणि}}
\textbf{\hi{मणि}} — \emph{maṇi} — \textquotedblleft{}a jewel; a gem\textquotedblright{}.

\begin{quote}
\hi{नाग} \hi{के} \hi{सिर} \hi{पर} \hi{एक} \hi{चमकदार} \hi{मणि} \hi{थी।}
\end{quote}

\begin{quote}
\emph{nāg ke sir par ek camakdār maṇi thī.}
\end{quote}

\begin{quote}
The cobra had a shining jewel on its head.
\end{quote}

Feminine.

\begin{cousinweb}
\textbf{\hi{मणि}} (\emph{maṇi}) is Sanskrit \emph{maṇi}, \textquotedblleft{}a jewel, a gem\textquotedblright{} (Monier-Williams). Old stories tell of a jewel in a serpent's hood.
\end{cousinweb}

\begin{grammarlens}[title={What you've built}]
One more word for telling what happened.
\end{grammarlens}

\subsection*{Your turn}
\begin{itemize}
\raggedright
  \item \emph{Say it:} \emph{maṇi}
  \item \emph{Say it:} \emph{maṇi}, once more
  \item \emph{Say it:} \emph{nāg ke sir par ek camakdār maṇi thī.} — The cobra had a shining jewel on its head.
  \item \emph{Recall:} say the Hindi for to bury, then the Hindi for a cobra, then say \emph{maṇi} again
\end{itemize}

\begin{tcolorbox}[breakable,colback=teal!4,colframe=teal!35!black,title={Before you move on}]
What does \hi{मणि} mean? (\textquotedblleft{}A jewel; a gem\textquotedblright{}.) Say it once more.
\end{tcolorbox}
\section[\texorpdfstring{\hi{डसना} (ḍasnā)}{डसना (ḍasnā)}]{\hi{डसना} (ḍasnā) — to bite (of a snake)}

\label{lesson:HI-C337-dasna}



\begin{quote}
Before the new one: say the Hindi for a cobra, then the Hindi for a jewel.
\end{quote}

\subsection*{You'll want to know: \hi{डसना}}
\textbf{\hi{डसना}} — \emph{ḍasnā} — \textquotedblleft{}to bite (of a snake)\textquotedblright{}.

\begin{quote}
\hi{नाग} \hi{ने} \hi{चोर} \hi{को} \hi{डसा।}
\end{quote}

\begin{quote}
\emph{nāg ne cor ko ḍasā.}
\end{quote}

\begin{quote}
The cobra bit the thief.
\end{quote}

Past, with \hi{ने} (he or she …): \hi{उसने} \hi{किसान} \hi{को} \hi{डसा}, \emph{usne kisān ko ḍasā}.

\begin{cousinweb}
\textbf{\hi{डसना}} (\emph{ḍasnā}) goes back to Sanskrit \emph{daśati}, \textquotedblleft{}bites, stings\textquotedblright{} (CDIAL 6230). It is the bite of a snake or a scorpion, not of a dog.
\end{cousinweb}

\begin{grammarlens}[title={What you've built}]
One more word for telling what happened.
\end{grammarlens}

\subsection*{Your turn}
\begin{itemize}
\raggedright
  \item \emph{Say it:} \emph{ḍasnā}
  \item \emph{Say it:} \emph{ḍasnā}, once more
  \item \emph{Say it:} \emph{nāg ne cor ko ḍasā.} — The cobra bit the thief.
  \item \emph{Recall:} say the Hindi for a cobra, then the Hindi for a jewel, then say \emph{ḍasnā} again
\end{itemize}

\begin{tcolorbox}[breakable,colback=teal!4,colframe=teal!35!black,title={Before you move on}]
What does \hi{डसना} mean? (\textquotedblleft{}To bite (of a snake)\textquotedblright{}.) Say it once more.
\end{tcolorbox}
